% Generated draft: requires mathematical and source review.
\documentclass[11pt]{article}
\usepackage[margin=25mm]{geometry}
\usepackage{fontspec}
\setmainfont{Latin Modern Roman}
\usepackage{amsmath,amssymb,mathtools,unicode-math}
\setmathfont{Latin Modern Math}
\usepackage{enumitem,needspace,xurl,hyperref}
\setlength{\parindent}{0pt}
\setlength{\parskip}{5pt}
\setlength{\emergencystretch}{3em}
\newcommand{\N}{\mathbb{N}}
\newcommand{\Z}{\mathbb{Z}}
\newcommand{\Q}{\mathbb{Q}}
\newcommand{\R}{\mathbb{R}}
\newcommand{\C}{\mathbb{C}}
\newcommand{\F}{\mathbb{F}}
\newcommand{\A}{\mathbb{A}}
\newcommand{\PP}{\mathbb{P}}
\newcommand{\E}{\mathbb{E}}
\newcommand{\eps}{\varepsilon}
\newcommand{\dd}{\,\mathrm{d}}
\begin{document}
\section*{UnsolvedMath ID: 30003612}
\textbf{Source problem code:} catalogue entry
\subsection{Definitions and mathematical statement}
Let \(X\) be a smooth uniruled projective variety, as in the supplied source. The base field and any characteristic restrictions are unspecified. The source also does not define a \emph{breaking map}, specify the allowed varieties serving as its domain, or state whether the definition depends on additional data on \(X\). Consequently, the following is a conditional formulation of the requested finiteness assertion, pending those definitions.

Write \(\mathcal B(X)\) for the intended class of breaking maps to \(X\). For maps \(f:Y\to X\) and \(f_i:Y_i\to X\), rational factorization through \(f_i\) means the existence of a rational map \(g:Y\dashrightarrow Y_i\) for which the composite is defined on a dense open subset of \(Y\) and
\[
 f=f_i\circ g
\]
there. Any requirements that \(g\) be dominant or that these maps be defined over a particular field remain unspecified by the record.

\textbf{Question.} For every smooth uniruled projective variety \(X\) in the intended setting, does there exist an integer \(r\geq 0\) and maps
\[
 f_i:Y_i\to X\quad (1\leq i\leq r),\qquad f_i\in\mathcal B(X),
\]
such that
\[
 \forall\,(f:Y\to X)\in\mathcal B(X),\quad
 \exists\,i\in\{1,\ldots,r\}\quad
 \exists\,g:Y\dashrightarrow Y_i
 \quad\text{with}\quad f=f_i\circ g
\]
as rational maps? The finite family may depend on \(X\). Without the missing definition of \(\mathcal B(X)\) and the intended geometric setting, this schema does not determine a fully specified mathematical conjecture.
\subsection{Short English statement}
The question asks whether each smooth uniruled projective variety admits finitely many breaking maps through which every breaking map to it factors rationally. The supplied record leaves the meaning of breaking map and the base field unspecified, so the statement requires clarification before mathematical review.
\subsection{Sources}
\begin{itemize}
\item Supplied local UnsolvedMath registry record, permanent ID \(30003612\), title \emph{Finite Factorization Bases for Breaking Maps}, source problem code \texttt{catalogue entry}, in \emph{Oberwolfach Reports Open Problems}. The record gives proposed year 2017 and names no proposer.
\item Supplied report citation: \emph{Algebraic Geometry: Birational Classification, Derived Categories, and Moduli Spaces} (2018), \emph{Oberwolfach Reports}, pp.~2703--2767, DOI \texttt{10.4171/owr/2017/45}. Supplied source URL and literature reference labeled ``Oberwolfach Report 45 2017'': \url{https://doi.org/10.4171/owr/2017/45}. The record describes these reports as documenting mathematical workshops and their open-problem discussions.
\item Supplied literature reference labeled ``Lehmann Tanimoto 2019'': \url{https://doi.org/10.1112/S0010437X19007103}. No article title or precise supporting location is supplied.
\item Literature assessment checked at \texttt{2026-08-21}: the supplied curation labels the problem open and describes the proposed finiteness as a central conjectural ingredient in the geometric Manin program with no general proof. This is supplied dated curation derived from the OWR record, not an independently verified literature assessment. The background expressly requests independent verification of time-sensitive status and cited claims.
\item The registry assigns the statement status \texttt{corrected\_verified}, explaining that mathematical content was preserved while notation, typography, or omitted definitions were repaired for self-contained publication. This is a supplied editorial label; it does not independently establish completeness or mathematical correctness of this draft.
\end{itemize}
\end{document}
